% ---
% title: ODE has no GRPO
% date: 2025/12/22
% tag: Diffusion
% ---
\section{问题}

为什么 deterministic flow 会导致计算 \(p(x_{t−1} | x_t, c)\) 变得“computationally expensive”，以及“divergence estimation”到底指什么?

\textbf{🔥 原句}

\begin{quote}
This deterministic approach fails to meet the GRPO policy update requirements in two key ways:
\end{quote}

\begin{quote}
\textbf{(1) rᵗᵢ(θ) in Eq.5 requires computing p(xₜ₋₁ | xₜ, c), which becomes computationally expensive under deterministic dynamics due to divergence estimation.}
\end{quote}

翻译：

\begin{quote}
该确定性（deterministic）的生成过程无法满足 GRPO 对策略更新的要求，其中一个关键原因是：
\end{quote}

\begin{quote}
GRPO 公式中的 rᵗᵢ(θ) 需要计算 \(p(xₜ₋₁ | xₜ, c)\)。
\end{quote}

\begin{quote}
但在确定性动力系统（ODE）下，这个条件分布的计算需要估计 divergence，从而非常耗费计算资源。
\end{quote}

\hrule

\section{🚧 为什么 deterministic flow 下无法直接计算 \(p(x_{t-1} | x_t, c)\)？}

因为 ODE dynamics 是 \textbf{完全确定的映射}：

\[
x_{t-1} = F_\theta(x_t, t)
\]

因此：

\[
p(x_{t-1}|x_t,c) = \delta(x_{t-1} - F_\theta(x_t,t))
\]

这是一个 \textbf{狄拉克 delta 分布}，它没有标准意义上的 log-probability。

但 GRPO（policy gradient）要求：

\[
\log \pi_\theta(a_t | s_t) = \log p_\theta(x_{t-1}|x_t,c)
\]

可是 δ 分布的 \(log\) 是 undefined：

\[
\log \delta(0) = \infty, \quad \log \delta(x\neq 0) = -\infty
\]

因此：

\begin{quote}
\textbf{ODE 模型没有可用的概率密度，无法用于 policy gradient。}
\end{quote}

\hrule

\section{🧠 进一步：为什么 deterministic dynamics 下“计算这个 density 会需要 divergence estimation”？}

这是这一句的核心。

上面说 δ 分布不可用，但严格来说：

你 \emph{可以} 用 \textbf{change-of-variables} 来定义 deterministic 的 density：

\[
p(x_{t-1}|x_t) = p(x_t|x_{t-1}) \cdot \left|\det\left(\frac{\partial x_t}{\partial x_{t-1}}\right)\right|^{-1}
\]

但在 flow matching（及 SD3 / Flux）中：

\begin{itemize}
  \item 模型从来没有训练过这个 Jacobian
  \item 也没有训练对数行列式（log-det-Jacobian）
\end{itemize}

要计算这个 density，你必须估计：

\[
\text{div}(v_\theta) = \nabla \cdot v_\theta
\]

或：

\[
\det\left(I + \Delta t \, \partial v_\theta / \partial x\right)
\]

这称为 \textbf{divergence estimation}。对 Transformer-based flow model 来说：

\textbf{❌ 完全不可能直接算}

Jacobian 是一个 4K × 4K 的矩阵（对 64×64×3 或更大图像来说是几十万维）

\textbf{❌ Hutchinson trick 也极其昂贵}

每次需要多次前向传播估计 trace。

\textbf{❌ 模型结构根本没有优化过 log-density}

Flow Matching 根本不是 normalizing flow，它不保证可逆性，也不保证 Jacobian tractable。

因此论文说：

\begin{quote}
computing \(p(x_{t−1} | x_t)\) becomes computationally expensive due to divergence estimation
\end{quote}

意思是：

如果你坚持要为 deterministic ODE 定义概率密度

\begin{itemize}
  \item 你必须计算 divergence（或 Jacobian determinant）
  \item 这是 \textbf{不可能也不必要}的操作
  \item GRPO 的 update 会完全崩溃
\end{itemize}

\hrule

简单说：\(p(x_{t-1} | x_t) = \delta(x_{t-1}-F_\theta(x_t))\) 不可求导，因此必须用雅可比divergence estimate，但是雅可比成本又太高了。

对于ODE而言，\( p(x_{t-1} | x_t) = \delta(x_{t-1}-F_\theta(x_t)) \)写作这个，意思是预测噪声刚好等于去噪前后的差的时候概率是无穷大（单点无穷），其他情况全是0，而且仍保持概率积分为1（只是定义成这样，表达出来他不可以求导）

唯一办法是使用连续可逆变换的定理：

\[
p(x_{t-1} | x_t) = p(x_t | x_{t-1}) \cdot \left|\det \frac{\partial x_t}{\partial x_{t-1}}\right|^{-1}
\]

其中核心项：

\[
\det \frac{\partial x_t}{\partial x_{t-1}}
\]

就是 \textbf{Jacobian determinant}。

而：

\[
\log \det (I + \Delta t \cdot \nabla_x v_\theta)
\]

展开后就是 \textbf{divergence（散度）相关项}：

\[
\text{div}(v_\theta) = \nabla\cdot v_\theta
\]

这就是所谓的 divergence estimation。

\hrule

下面我将 \textbf{非常清晰、数学严格、一步不跳地推导}：

\begin{quote}
\textbf{为什么 deterministic flow 的概率密度会出现}
\end{quote}

\begin{quote}
\(\log \det (I + \Delta t \cdot \nabla_x v_\theta)\)
\end{quote}

\begin{quote}
这个项到底从哪里来的？
\end{quote}

这个式子其实来自： ⭐ 离散化的 ODE → 一个可逆映射的 Jacobian determinant

\hrule

\section{🔥 核心背景：我们想给 deterministic 生成过程定义概率密度}

Flow Matching 的生成过程（反向）本质是一个 ODE：

\(\frac{dx}{dt} = v_\theta(x,t)\)

如果使用 Euler 法离散化：

\[
x_{t-\Delta t} = x_t - v_\theta(x_t,t)\Delta t
\]

这是一个确定性映射：

\[
x_{t-1} = F(x_t)
\]

我们想计算\textbf{这个映射的概率密度变化}，使其能被 policy gradient 使用。

概率密度变化由 change-of-variables 定律决定：

\section{🧠 Change-of-variables formula}

对一个可逆映射：\(x' = F(x)\)

概率密度满足：

\[
p_{x'}(x') = p_x(x)\cdot \left|\det \frac{\partial F(x)}{\partial x}\right|^{-1}
\]

其中：

\[
\frac{\partial F(x)}{\partial x}
\]

就是 \textbf{Jacobian matrix}。

\section{🔍 应用到 flow matching 的离散 ODE}

映射：

\[
x_{t-1} = F(x_t) = x_t - v_\theta(x_t,t)\Delta t
\]

对 \(x_t\) 求 Jacobian：

\[
J = \frac{\partial x_{t-1}}{\partial x_t}
\]

展开：

\[
x_{t-1} = x_t - \Delta t\cdot v_\theta(x_t,t)
\]

求导：

\[
\frac{\partial x_{t-1}}{\partial x_t} = I - \Delta t \cdot \frac{\partial v_\theta}{\partial x_t}
\]

也就是：

\[
J = I - \Delta t\,\nabla_x v_\theta
\]

因此：

\[
\log \left|\det J\right| = \log \left|\det\left(I - \Delta t \cdot \nabla_x v_\theta\right)\right|
\]

\section{📌 这个 Jacobian determinant 和 divergence 的关系是什么？}

对于小步长 Δt（Euler 离散），使用一阶泰勒展开：

\[
\log \det(I + \epsilon A) \approx \epsilon \, \text{tr}(A)
\]

于是：

\[
\log \det (I + \Delta t \cdot \nabla_x v_\theta) \approx \Delta t\cdot \text{tr}(\nabla_x v_\theta)
\]

而：

\[
\text{tr}(\nabla_x v_\theta) = \nabla \cdot v_\theta
\]

就是 \textbf{divergence（散度）}。

\section{🔥 最终结论（论文想表达的真正意思）}

为了计算 deterministic flow 的概率密度，你需要：

\begin{enumerate}
  \item 计算 Jacobian：
  
  \[
  I + \Delta t\cdot\nabla_x v_\theta
  \]
  \item 计算它的 determinant：
  
  \[
  \det(I + \Delta t\nabla_x v_\theta)
  \]
  \item 计算它的 log：
  
  \[
  \log \det(I + \Delta t\nabla_x v_\theta)
  \]
\end{enumerate}

但这在图像生成模型中意味着：

\begin{itemize}
  \item Jacobian 是上万×上万的矩阵
  \item det 的计算是  \(O(n^3)\)
  \item divergence 用 Hutchinson trick 也需要多次前向传播
\end{itemize}

\textbf{完全不可行。}

所以 deterministic flow \textbf{不能直接做 GRPO}。

这就是论文强调的：

\begin{quote}
computing \(p(x_{t−1}|x_t)\) is expensive due to divergence estimation.
\end{quote}

\section{✨ 一句话总结}

\[
\log \det(I + \Delta t \nabla_x v_\theta)
\]

来源于：

\begin{itemize}
  \item 对 flow matching 的 ODE 离散化
  \item 应用 change-of-variables formula
  \item 得到概率密度更新项
  \item 其一阶近似就是 divergence
\end{itemize}

这项计算量巨大，使 deterministic flow 无法用于 RL。
